\subsection{No-Deleting Principle}\label{sec:no-deleting-principle}

The \textbf{No-Deleting Principle} is the dual counterpart of the no-cloning theorem. While the no-cloning theorem states that it is impossible to perfectly copy an arbitrary unknown quantum state, the no-deleting principle states that it is also \textbf{impossible to perfectly delete quantum information through a unitary operation}.

\highspace
In other words:
\begin{itemize}
    \item Quantum information \textbf{cannot} be freely \textbf{copied};
    \item Quantum information \textbf{cannot} be freely \textbf{destroyed}.
\end{itemize}
Like the no-cloning theorem, the no-deleting principle is \hl{deeply connected to the unitary structure} of quantum mechanics.

\highspace
\begin{definitionbox}[: No-Deleting Principle]
    Suppose we want a quantum gate $U$ that deletes a qubit state:
    \begin{equation*}
        U\ket{x}=\ket{0}
    \end{equation*}
    For an arbitrary unknown state $\ket{x}$.

    \highspace
    Equivalently, if we have two identical copies of a state, we may wish to delete one copy while leaving the other unchanged:
    \begin{equation}
        U\ket{x}\ket{x} = \ket{x}\ket{0}
    \end{equation}
    The \definition{No-Deleting Theorem} states that \hl{no universal unitary gate} capable of performing this operation \hl{can exist} for arbitrary unknown quantum states.

    \highspace
    Intuitively, this means that \textbf{quantum information cannot} simply \textbf{disappear through reversible quantum evolution}.
\end{definitionbox}

\highspace
As with the no-cloning theorem, we can prove the no-deleting principle by contradiction, using the fact that unitary operators preserve inner products (see \autopageref{deep-box:alternative-proof-no-cloning}).

\begin{proof}
    Suppose, by contradiction, that there exists a deleting gate $U$ such that:
    \begin{equation*}
        U\ket{x}=\ket{0}
    \end{equation*}
    For an arbitrary quantum state $\ket{x}$.

    \highspace
    Since $U$ is supposed to be universal, it must also delete another arbitrary state $\ket{y}$:
    \begin{equation*}
        U\ket{y}=\ket{0}
    \end{equation*}
    At this point, we have two equations describing the action of $U$ on two different input states.

    \highspace
    Now, the only thing we know about unitary operators is that they \textbf{preserve inner products}. Therefore, if $U$ is unitary, the inner product between the input states must be equal to the inner product between the corresponding output states.

    \begin{itemize}
        \item \textbf{Before deleting}, the inner product between the two input states is:
        \begin{equation*}
            \left(U\ket{x}\right) \cdot \left(U\ket{y}\right)
        \end{equation*}

        Using the definition of inner product:
        \begin{align*}
            \left(U\ket{x}\right)\cdot\left(U\ket{y}\right)
            &=
            \left(U\ket{x}\right)^{\dagger} \left(U\ket{y}\right) \\[.2em]
            %
            &=
            \bra{x} U^{\dagger} U \ket{y} \\[.2em]
            %
            &\downarrow \text{since $U$ is unitary, $U^{\dagger} U = I$} \\[.2em]
            %
            &=
            \braket{x}{y}
        \end{align*}

        \item \textbf{After deleting}, both output states become $\ket{0}$, so the inner product is:
        \begin{equation*}
            \ket{0}\cdot\ket{0}
        \end{equation*}

        Again:
        \begin{align*}
            \ket{0}\cdot\ket{0}
            &=
            \ket{0}^{\dagger}\ket{0} \\[.2em]
            %
            &=
            \braket{0}{0} \\[.2em]
            %
            &\downarrow \text{since $\ket{0}$ is a normalized state (\autopageref{def:normalized-quantum-state})} \\[.2em]
            %
            &=
            1
        \end{align*}
    \end{itemize}
    Since unitary operators preserve inner products:
    \begin{equation*}
        \braket{x}{y}=1
    \end{equation*}
    But this is \textbf{not true in general for arbitrary quantum states}. It is true only if:
    \begin{equation*}
        \ket{x}=\ket{y}
    \end{equation*}
    Therefore, such a \hl{universal deleting gate $U$ cannot exist.}
\end{proof}